\subsection{Harman 421} % todo: verwijder deze heading en laat enkel de tekst staan

According to the International Institute for Strategic Studies (IIIS), in 2025, global defense spending reached USD2.63 trillion, with the United States, China and Russia being the front runners~\citep{iiss2026militarybalance}.
Raw expenditure figures, however, show an incomplete picture. The three aforementioned countries have very large populations, with China sitting at an estimated population of 1.4 billion, the United States at 330 million and Russia at 144 million. When defense spending is expressed as a percentage of gross domestic product (GDP), we get a more accurate overview of the biggest spenders.

By this metric Ukraine spends 34\% of their GDP on military expenditure~\citep{sipri2025milex}. This is due to Ukraine having been in conflict with Russia since 2014, with the conflict escalating into a full-blown war since the 2022 Russian invasion of Ukraine. Second is Israel, which spends 8.8\% of its GDP on defense~\citep{sipri2025milex}. This number is not surprising given the geopolitical situation in the Middle East; Since its inception in 1948, Israel has almost constantly been in conflict with its neighbors, fighting at least five to six major Arab-Israeli wars, along with many smaller conflicts.

This large expendature is reflected in the country's use of artificial intelligence (AI)-assisted autonomous weapon systems. Israel's use of artificial intelligence in warfare is very well documented, with the country having developed many of its own AI systems. AI-driven systems like `Lavender' and `Habsora' are used to autonomously identify bombing targets, with `Lavender' being used to identify human targets and `Habsora' being used to identify military targets such as buildings and structures instead of individuals~\citep{abraham2024lavender}.

Lavender and Habsora are the focus of this report. Most recently, they were used in the 2023 Israel-Hamas war, where they generated target lists for the Israeli Air Force (IAF)~\citep{abraham2024lavender}. Sources interviewed by +972 Magazine and Local Call claim that \textquote{during the first weeks of the war, the army almost completely relied on Lavender, which clocked as many as 37,000 Palestinians as suspected militants \textemdash{} and their homes \textemdash{} for possible air strikes.}~\citep{abraham2024lavender}. A problem that arises in this scenario is that active military personnel start working on a "rubber stamp" basis, where they are not verifying the target lists generated by the AI system and are instead completely relying on the AI system to make decisions that have life and death consequences. One source admitted that human overseers only devoted 20 seconds to each target, were merely checking if the target was male, before approving a strike. This effectively delegates life-and-death decisions to an AI system, which makes establishing accountability and responsibility rather difficult.